\section*{Bab \ref{ch:b3:planetary-atmospheres} --- Atmosfer Planet dan Kelayakhunian}

\begin{solution}{exo:b3:planetary-atmospheres:1}
$H = RT/Mg$: Venus $8.314\times737/(0.04401\times8.87) = \qty{15.7}{km}$;
Bumi $8.314\times288/(0.02897\times9.80) = \qty{8.4}{km}$; Mars
$8.314\times214/(0.04334\times3.71) = \qty{11.1}{km}$; Titan
$8.314\times94/(0.0280\times1.35) = \qty{20.7}{km}$. Atmosfer Titan,
meskipun sedingin itu, adalah yang paling luas: gravitasinya tujuh kali lebih lemah
daripada gravitasi Bumi.
\end{solution}

\begin{solution}{exo:b3:planetary-atmospheres:2}
$\lambda = GMm/kTr = \num{6.674e-11}\times\num{5.97e24}\times\num{1.674e-27}
/(\num{1.381e-23}\,T\times\num{6.878e6})$: 8.78 pada \qty{800}{K}, 4.39 pada
\qty{1600}{K}. Fraksi yang lolos adalah $9.78\,e^{-8.78} =
\num{1.5e-3}$ dan $5.39\,e^{-4.39} = \num{6.7e-2}$: siklus Matahari
mengubah fraksi lolos hidrogen dengan faktor sekitar 44.
\end{solution}

\begin{solution}{exo:b3:planetary-atmospheres:3}
(a)~$d = S_\mathrm{eff}^{-1/2}$: \omterm{def:b3:planetary-atmospheres:runaway}{rumah kaca lembap} \qty{0.993}{au}, \omterm{def:b3:planetary-atmospheres:runaway}{rumah kaca tak terkendali}
\qty{0.975}{au}, rumah kaca maksimum \qty{1.705}{au}. (b)~Venus, pada 1.91,
menerima lebih dari fluks \omterm{def:b3:planetary-atmospheres:runaway}{batas rumah kaca tak terkendali} dan bahkan lebih dari batas empiris
Venus terkini (1.78): ia berada di dalam tepi dalam. Mars, pada 0.43,
berada antara 1.014 dan 0.3438, di dalam zona layak huni konservatif:
kekeringannya berasal dari ukurannya yang kecil dan hilangnya atmosfernya, bukan
dari jaraknya.
\end{solution}

\begin{solution}{exo:b3:planetary-atmospheres:4}
(a)~$\log_{10}p = 3.9895 - 443.028/93.66 = -0.741$, $p =
\qty{0.18}{bar}$. (b)~$0.05\times1.6 = \qty{0.080}{bar}$, kelembapan
relatif sebesar $0.080/0.18 = 44\,\%$. (c)~Udara di atas lautan jauh dari
jenuh: metana menguap dari lautan itu, seperti air menguap dari
samudra Bumi, dan udara lembap harus naik dan mendingin sebelum hujan turun.
\end{solution}

\begin{solution}{exo:b3:planetary-atmospheres:5}
$\lambda_\mathrm D = 7.02\times2.014/1.008 = 14.0$. $u_\mathrm D/u_\mathrm H
= \sqrt{1.008/2.014} = 0.71$. Rasio per atom $= 0.71\times(15.0\,
e^{-14.0})/(8.02\,e^{-7.02}) = 0.71\times\num{1.69e-3} = \num{1.2e-3}$.
\omterm{def:b3:planetary-atmospheres:jeans-escape}{Lolos Jeans} hampir tidak melepaskan deuterium: ia penyaring kuat, dengan
$\alpha \approx 0.001$.
\end{solution}

\begin{solution}{exo:b3:planetary-atmospheres:6}
(a)~$0.016/\num{1.5575e-4} = 103$. (b)~Dengan $\alpha = 0$, pengayaannya adalah
$M_0/M$: $M_0 = 103\times\num{6e15} = \qty{6.2e17}{kg}$, atau
$\num{6.2e17}/(1000\times\num{4.60e14}) = \qty{1.3}{m}$ lapisan air di seluruh
Venus. (c)~$M_0/M = 103^{1/0.2} = \num{1.2e10}$, persediaan masa lalu sebesar
\qty{7e25}{kg}, dua belas kali massa Bumi: mustahil. Pengayaan
memerlukan lolos berfraksionasi kuat (\omterm{def:b3:planetary-atmospheres:jeans-escape}{lolos Jeans} atau \omterm{def:b3:planetary-atmospheres:hydrodynamic}{lolos nontermal}),
aktif setidaknya selama hilangnya sisa air terakhir, atau keadaan
tunak dengan pasokan air.
\end{solution}

\begin{solution}{exo:b3:planetary-atmospheres:7}
(a)~Dalam empat miliar tahun ($\num{1.26e17}$~s) Mars kehilangan
$\num{3.8e43}$ atom, $\num{1.9e43}$ molekul air, $\qty{5.7e17}{kg}$;
pada luas $4\pi(\num{3.3895e6})^2 = \qty{1.44e14}{m^2}$, yaitu lapisan
setebal \qty{3.9}{m}. (b)~Dua puluh kali lebih banyak, \qty{78}{m}. (c)~Bahkan laju
pada kemiringan tinggi hanya mencapai ujung bawah rentang 100 sampai
\qty{1500}{m}; sisanya masuk ke kerak dan es, atau hilang
lebih cepat ketika Matahari masih muda.
\end{solution}

\begin{solution}{exo:b3:planetary-atmospheres:8}
(a)~$m = pA/g = 636\times\num{1.44e14}/3.71 = \qty{2.5e16}{kg}$.
(b)~$0.1\times\num{1.26e17} = \qty{1.3e16}{kg}$, setengah
atmosfer sekarang; seluruhnya dalam $\num{2.5e16}/0.1 = \qty{2.5e17}{s}$,
\num{7.8}~miliar tahun. (c)~Matahari muda lebih aktif dan badainya
lebih sering, dan kehilangan membesar selama badai; saluran lain
(hidrogen, sputtering, lolos fotokimia) tidak dihitung.
\end{solution}

\begin{solution}{exo:b3:planetary-atmospheres:9}
(a)~$S(1 - 0.193)/4 = 291$ memberi $S = \qty{1442}{W.m^{-2}}$, 1.06 kali
Bumi. (b)~Venus hari ini akan menyerap $0.807\times2601/4 =
\qty{525}{W.m^{-2}}$; Venus muda, $0.7\times2601 = \qty{1821}{W.m^{-2}}$
diterima, \qty{367}{W.m^{-2}} diserap: keduanya di atas batas, sehingga
Venus muda tanpa awan sudah mengalami \omterm{def:b3:planetary-atmospheres:runaway}{rumah kaca tak terkendali} sejak awal. (c)~Awan
menaikkan albedo; model tiga dimensi dengan awan menemukan jalan
menuju Venus muda yang layak huni.
\end{solution}

\begin{solution}{exo:b3:planetary-atmospheres:10}
(a)~$H = kT/mg = \num{1.381e-23}\times1000/(16\times\num{1.661e-27}\times9.21)
= \qty{56}{km}$. (b)~$n_\mathrm{exo} = 1/(\sigma_c H) = 1/(\num{3e-19}\times
\num{5.64e4}) = \qty{5.9e13}{m^{-3}}$; $z = 200 + 56.4\ln(\num{1e16}/
\num{5.9e13}) = 200 + 56.4\times5.13 = \qty{490}{km}$. (c)~Gravitasi yang lebih lemah
di tempat lebih tinggi membuat tinggi skala membesar seiring ketinggian, sehingga massa jenis
turun lebih lambat daripada yang diasumsikan: \omterm{def:b3:planetary-atmospheres:exobase}{eksobase} naik.
\end{solution}

\begin{solution}{exo:b3:planetary-atmospheres:11}
(a)~$\Phi = \num{3.0e13}\times\num{3e-3} = \qty{9e10}{cm^{-2}.s^{-1}} =
\qty{9e14}{m^{-2}.s^{-1}}$; pada luas $4\pi(\num{6.371e6})^2 =
\qty{5.10e14}{m^2}$, $\num{4.6e29}$ molekul $\ce{H2}$ per detik.
(b)~Sebuah samudra menampung $\num{1.4e24}/18.015\times\num{6.022e23} =
\num{4.7e46}$ molekul air, masing-masing satu $\ce{H2}$: ia bertahan
$\qty{1.0e17}{s} = \num{3.2}\text{ miliar tahun}$, sebanding dengan umur
Bumi, itulah sebabnya rasio campuran \num{3e-3} menandai batas \omterm{def:b3:planetary-atmospheres:runaway}{rumah kaca
lembap}.
\end{solution}

\begin{solution}{exo:b3:planetary-atmospheres:12}
(a)~$p_1 = g/D\kappa = 9.81/(1.66\times0.005) = \qty{1180}{Pa}$.
(b)~$1/T = 1/273.16 - \ln(1182/611.657)\times461.5/\num{2.5e6}$, $T =
\qty{282.5}{K}$, dan $\sigma T^4 = \qty{361}{W.m^{-2}}$. (c)~Batas sebenarnya
lebih rendah, sehingga level radiasi lebih dingin dan pada tekanan lebih rendah: uap air
sebenarnya lebih buram daripada gas kelabu ini, dan $\kappa$ efektifnya
lebih besar.
\end{solution}

\begin{solution}{pb:b3:planetary-atmospheres:1}
\textbf{1.} $g = GM/R^2 = \num{6.674e-11}\times\num{5.2e24}/(\num{6.2e6})^2 =
\qty{9.03}{m.s^{-2}}$; $\bar\rho = M/\tfrac43\pi R^3 = \qty{5210}{kg.m^{-3}}$.

\textbf{2.} $v_\mathrm{esc} = \sqrt{2GM/r} = \sqrt{2\times\num{3.47e14}/\num{6.7e6}}
= \qty{10.18}{km.s^{-1}}$.

\textbf{3.} $u = \sqrt{2kT/m}$: hidrogen \qty{4.81}{km.s^{-1}}, helium
\qty{2.41}{km.s^{-1}}, oksigen \qty{1.21}{km.s^{-1}}.

\textbf{4.} $\lambda = (v_\mathrm{esc}/u)^2$: 4.49, 17.8, 71.2.

\textbf{5.} Hanya hidrogen: fraksi lolosnya adalah 0.06, untuk helium
\num{3.5e-7}, untuk oksigen \num{9e-30}. Tidak satu pun dekat dengan rezim
hidrodinamik, $\lambda \approx 1.5$ hingga 3.

\textbf{6.} $\lambda = 4.49\times1400/3000 = 2.09$: berada dalam transisi;
hidrogen akan sudah mengalir lepas secara hidrodinamik, mungkin sambil menyeret
gas-gas yang lebih berat.

\textbf{7.} $(1 + 4.49)\,e^{-4.49} = 0.0618$.

\textbf{8.} $n u/2\sqrt\pi = \num{2e11}\times4806/3.545 = \qty{2.71e14}{m^{-2}.s^{-1}}$;
$\Phi_\mathrm J = \num{2.71e14}\times0.0618 = \qty{1.68e13}{m^{-2}.s^{-1}}$.

\textbf{9.} $4\pi(\num{6.7e6})^2 = \qty{5.64e14}{m^2}$; $\num{9.5e27}$ atom per
detik.

\textbf{10.} $\num{1.4e24}/18.015 = \qty{7.8e22}{mol}$, $\num{4.68e46}$ molekul,
$\num{9.36e46}$ atom hidrogen.

\textbf{11.} $\num{9.36e46}/\num{9.5e27} = \qty{9.9e18}{s}$, sekitar 310 miliar
tahun.

\textbf{12.} $\lambda_\mathrm D = 4.49\times2.014/1.008 = 8.96$; fraksi lolos
$9.96\,e^{-8.96} = \num{1.28e-3}$; efisiensi relatif $0.707\times
\num{1.28e-3}/0.0618 = 0.015$.

\textbf{13.} $H_a = kT/m_a g = \num{1.381e-23}\times200/(44.01\times\num{1.661e-27}
\times9.03) = \qty{4.18}{km}$.

\textbf{14.} $b = \qty{1.2e21}{m^{-1}.s^{-1}}$, sehingga
\[
\Phi_\mathrm{DL} = \frac{b f(m_a - m_i)}{m_a H_a}
= \frac{\num{1.2e21}\times\num{1e-3}\times(42.0/44.0)}{4185}
= \qty{2.74e14}{m^{-2}.s^{-1}}.
\]

\textbf{15.} $\num{3.0e13}\times\num{1e-3}\times9.03/9.81 =
\qty{2.76e10}{cm^{-2}.s^{-1}} = \qty{2.76e14}{m^{-2}.s^{-1}}$: sama
dalam 1\,\%.

\textbf{16.} $4\pi(\num{6.2e6})^2 = \qty{4.83e14}{m^2}$; $\num{1.32e29}$ molekul
dari $\ce{H2}$, $\num{2.64e29}$ atom per detik.

\textbf{17.} Batas difusi menyediakan 28 kali lipat dari yang dilepaskan \omterm{def:b3:planetary-atmospheres:exobase}{eksobase} pada
massa jenisnya saat ini. Hidrogen menumpuk di termosfer sampai
fluks Jeans menyamai pasokannya; dalam jangka panjang lolosnya
terbatas-difusi, seperti di Bumi.

\textbf{18.} $\num{2e11}\times28 = \qty{5.6e12}{m^{-3}}$.

\textbf{19.} $\num{4.68e46}/\num{1.32e29} = \qty{3.54e17}{s}$, 11.2 miliar tahun.

\textbf{20.} $S = 0.30/0.60^2 = 0.833$.

\textbf{21.} $d = \sqrt{L/S_\mathrm{eff}}$: tepi dalam $\sqrt{0.30/0.939} =
\qty{0.565}{au}$, tepi luar $\sqrt{0.30/0.290} = \qty{1.02}{au}$. Kepos,
pada \qty{0.60}{au}, berada di dalam, dekat dengan tepi dalam.

\textbf{22.} $0.939/0.833 = 1.13$: peningkatan kecerahan sebesar 13\,\%.

\textbf{23.} $0.833\times1361\times0.7/4 = \qty{198}{W.m^{-2}}$, $T_\mathrm{eq} =
(198/\sigma)^{1/4} = \qty{243}{K}$; jauh di bawah
\qty{291}{W.m^{-2}} sebagai batas (dihitung untuk bintang mirip Matahari, tetapi
selisihnya cukup lebar).

\textbf{24.} Tiga kali lebih cepat: 3.7 miliar tahun.

\textbf{25.} \textbf{11~miliar~tahun}.
\end{solution}
